\section{Flavor parameters, loop removal, and stability}
\label{sec:deformations}

Throughout this section, $B_F$ denotes the quotient by the full topological
torus.  In particular, it is independent of the stability character.
We first prove that regular stabilities have the same Betti numbers.
We then lift a no-flavor comparison to arbitrary flavor parameters and
use an added-loop parameter to remove the loop hypothesis at determinant
stability.  These three inputs combine in Theorem~\ref{thm:full-torus}.

\begin{lemma}\label{lem:betti-independence}
Fix a quiver and dimension and framing vectors.  The ordinary Betti numbers
of $X_\theta=\mu^{-1}(0)/\!/_{\theta}G$ are independent of the regular
character $\theta$, where regular means that every semistable point is stable
and the action on the semistable locus is free.  Empty quotients are included.
\end{lemma}

\begin{proof}
By Lemma~\ref{lem:local-stability}, each regular character has a real open
ball in which every rational parameter gives the same stable and semistable
loci, geometric quotient, and principal gauge bundle.  We may therefore
perturb within this ball without changing the quotient.

Here is a comparison of such neighborhoods for different regular
characters.  Write $G_c=\prod_i U(V_i)$ for a compact form and
$\mathfrak c=(\mathfrak g^*)^G$ for the space of complex central moment
parameters.  Equip $T^*N$ with its standard flat hyperk\"ahler structure.
The character space identifies with a real form of $\mathfrak c$.
For a rational parameter, clear its denominators and then dilate the
flat representation by the inverse square root of that common denominator;
the quadratic moment maps thereby return to level $\theta$.
Kempf--Ness identifies $X_\theta$ with the quotient by $G_c$ at real moment
level $\theta$ and complex moment level zero.  Quaternionic rotation commutes
with $G_c$ and rotates the three moment maps.  Thus it identifies this space
with the quotient at real level zero and complex level $c\theta$, where
$c\ne0$ is one fixed scalar determined by the normalization of the moment
maps.  Applying affine Kempf--Ness at this latter level gives a homeomorphism
\begin{equation}\label{eq:rotation-affine-fiber}
 X_\theta\ \simeq\ \mu^{-1}(c\theta)/\!/G .
\end{equation}
This uses the flat hyperk\"ahler quotient construction
\cite[Section~3(D)]{HKLR87} and the topological affine Kempf--Ness
identification recalled in Section~\ref{sec:preliminaries}; the
quiver-variety description appears in
\cite{Nakajima94}.  The rotation here acts directly on the flat
representation and is valid with loops.  It requires no properness of a
moment-map family.  One can normalize $c=1$; alternatively a linear dilation
of $T^*N$ identifies any two common nonzero scalar multiples of the complex
level.

Consider the affine family
\[
 \pi:\mathcal Z=
 \Spec\C[\{(x,\zeta)\in T^*N\times\mathfrak c:
                       \mu(x)=\zeta\}]^G
 \longrightarrow\mathfrak c.
\]
Its fiber at $\zeta$ is $\mu^{-1}(\zeta)/\!/G$.
To check this without a flatness assumption, let $A$ be the coordinate ring
inside the brackets.  The ideal specifying $\zeta$ is generated by invariant
parameter functions.  Exactness of invariants and the Reynolds operator
give
\[
 (A/\mathfrak m_\zeta A)^G
       =A^G/\mathfrak m_\zeta A^G.
\]
The morphism $\pi$ is of finite type.  Constructibility of compactly
supported direct image and its fiber-stalk description
\cite[Theorem~2.5(b)]{MaximSchurmann21}
\cite[Theorem~5.10]{SchnurerSoergel16}
therefore imply that $R\pi_!\C_{\mathcal Z}$ has constructible cohomology
sheaves with stalks $H_c^*(\pi^{-1}(\zeta);\C)$.  On a common
dense Zariski open subset $\Omega\subset\mathfrak c$, their ranks, hence all
compactly supported Betti numbers of the fibers, are constant.  We may take
$\Omega$ connected.  Every real open ball in the character space is Zariski
dense after multiplication by $c$, and its rational points are dense in the
usual topology.  Each regular character can thus be perturbed, without
changing its quotient, so that $c\theta\in\Omega$.

By Equation~\eqref{eq:rotation-affine-fiber}, all regular quotients have the same
compactly supported Betti numbers.  A nonempty such quotient is smooth of
complex dimension $2(\dim N-\dim G)$: freeness makes the moment map a
submersion along its stable locus.  Poincar\'e duality now gives equality of
ordinary Betti numbers.  If one regular quotient is empty, all compactly
supported Betti numbers vanish for every regular quotient; a nonempty smooth
quotient would have nonzero top compactly supported cohomology, a
contradiction.  This also covers nonfaithful gauge actions and unframed
components: the generic central fiber is allowed to be empty.  If $G$ is
trivial, the parameter space is a point and the assertion is immediate.
\end{proof}

\begin{lemma}\label{lem:flavor-lift}
For a fixed regular stability, suppose that the coefficient map induces
$B_0\simeq H^*(X)$.  Then, for every commuting flavor torus $F$, the map of
Proposition~\ref{prop:kernel-inclusion} is an isomorphism
\[
 B_F\xrightarrow{\ \sim\ }H_F^*(X)
\]
of graded $R$-algebras.
\end{lemma}

\begin{proof}
Set $S=\C[\mathfrak f]$ and let $S_+$ be its positive-degree ideal.
Kirwan surjectivity implies $H^{\mathrm{odd}}(X)=0$.
The Borel spectral sequence for $X\to EF\times_F X\to BF$ has second page
$H^*(BF)\otimes H^*(X)$, concentrated in even total degrees.  It degenerates,
and its edge map to $H^*(X)$ is surjective.  Lifting a finite homogeneous
basis of $H^*(X)$ and applying Leray--Hirsch shows that $H_F^*(X)$ is a free
$S$-module with
\[
 H_F^*(X)/S_+H_F^*(X)=H^*(X).
\]
This remains valid for $X=\varnothing$, with zero modules.

Let $J$ be the kernel of the surjection
$B_F\twoheadrightarrow H_F^*(X)$ from
Proposition~\ref{prop:kernel-inclusion}.  Freeness of the target makes
$\operatorname{Tor}_1^S(H_F^*(X),S/S_+)$ vanish.  Tensoring its defining exact
sequence with $S/S_+$ consequently gives
\[
 0\longrightarrow J/S_+J\longrightarrow B_F/S_+B_F
       \longrightarrow H^*(X)\longrightarrow0.
\]
Lemma~\ref{lem:flavor-base-change} identifies the middle term and map with
$B_0\xrightarrow{\sim}H^*(X)$; hence $J=S_+J$.
The module $J$ is graded and bounded below, since $B_F$ is a quotient of the
nonnegatively graded ring $R$.  If $J$ had a nonzero homogeneous element of
least degree, the identity $J=S_+J$ would express it using elements of smaller
degree, which is impossible.  Thus $J=0$.
This argument uses positivity only for the coefficient quotient $B_F$, and
does not assume that the Coulomb algebra itself is nonnegatively graded.
\end{proof}

\begin{proposition}\label{prop:loop-removal}
For every finite quiver and every dimension and framing vector, at the
determinant stability used in Corollary~\ref{cor:loop-isomorphism}, the
coefficient map induces an isomorphism of graded $R_0$-algebras
\[
 B_0\xrightarrow{\ \sim\ }H^*(X).
\]
\end{proposition}

\begin{proof}
Discard zero-rank gauge factors.  Add one loop at each remaining vertex and
write
\[
 N^+=N\oplus\bigoplus_i\End(V_i).
\]
Let $H=\C^\times$ act trivially on $N$ and with weight one on every added
loop.  Put $H^*_H(\mathrm{pt})=\C[z]$, with $\deg z=2$, and use a superscript
$+$ for the enlarged theory.  At determinant stability,
Corollary~\ref{cor:loop-isomorphism} and Lemma~\ref{lem:flavor-lift} give
\begin{equation}\label{eq:enlarged-flavor-isomorphism}
 B_H^+\xrightarrow{\ \sim\ }H_H^*(X^+)
 \qquad\text{as graded }R_0[z]\text{-algebras}.
\end{equation}
We compare this with the original theory endowed with the trivial $H$-action.
Let $\mathfrak m_0\subset R_0$ be the maximal ideal of the graded origin and
set
\begin{equation}\label{eq:loop-local-ring}
 \Lambda=R_0[z]_{(\mathfrak m_0,z-1)}.
\end{equation}
The comparison over $\Lambda$ is ungraded in cohomological degree but
retains topological charge.  The equality of homogeneous ideals in $R_0$
obtained below will restore the cohomological grading.

\emph{The Coulomb comparison.}
Use the common rational Laurent realization of
Proposition~\ref{prop:minuscule}, and put $P=\C[x_{ia}]$.
For a positive or negative minuscule coweight $\lambda$, the added numerator
in Equation~\eqref{eq:minuscule} is
\[
 E_\lambda(x,z)=
   \prod_i\prod_{a,b=1}^{v_i}
     (z+x_{ib}-x_{ia})^{(\lambda_{ia}-\lambda_{ib})_+}.
\]
It is invariant under the stabilizer $\mathcal W_\lambda$ and has value one
at $(x,z)=(0,1)$.  Its inverse belongs to
$P[z]^{\mathcal W_\lambda}\otimes_{R_0[z]}\Lambda$: indeed,
\[
 \Delta_\lambda=\prod_{w\in\mathcal W}wE_\lambda
       \in R_0[z],\qquad \Delta_\lambda(0,1)=1,
\]
and $E_\lambda^{-1}=(\Delta_\lambda/E_\lambda)/\Delta_\lambda$.
If $M(\lambda,p)$ denotes an original dressed generator, the enlarged
generator with dressing $p$ is $M(\lambda,pE_\lambda)$.  Conversely,
dressing by $p/E_\lambda$ recovers $M(\lambda,p)$ after localization.
Generation therefore identifies the two localized Coulomb algebras, with
their coefficient maps and their full topological gradings.  The elements
inverted in Equation~\eqref{eq:loop-local-ring} have topological degree zero.  Thus
localization commutes with quotienting by all nonzero topological degrees,
and gives
\begin{equation}\label{eq:loop-coulomb-localization}
 B_H^+\otimes_{R_0[z]}\Lambda
    \simeq (B_0\otimes_\C\C[z])\otimes_{R_0[z]}\Lambda.
\end{equation}

\emph{The cohomological comparison.}
Let $D^+\subset T^*N^+$ be the stable zero moment locus.  An $H$-fixed
point of the quotient may be fixed only after a compensating gauge
transformation.  We therefore localize the actual $H$-fixed locus of this
prequotient, retaining $G\times H$ equivariance, and first restrict to
$T\times H$.  Brion's subgroup localization theorem
\cite[Lemma~4]{Brion98} says that, for a torus space embedded equivariantly in
a finite-dimensional representation, restriction to the fixed locus of a
subtorus becomes an isomorphism after inverting finitely many characters
nontrivial on that subtorus.  Apply it to the compact forms of $T\times H$
and $H$ and to the invariant locally closed subset $D^+\subset T^*N^+$.
Every character nontrivial on $H$ is nonzero at the central Lie algebra
element $(0,1)$.  The restriction
\[
 H^*_{T\times H}(D^+)\longrightarrow H^*_{T\times H}((D^+)^H)
\]
is consequently an isomorphism there.  Taking $\mathcal W$-invariants
recovers $G\times H$ equivariance.  This is compatible with localization
over $\Lambda$: the point is $\mathcal W$-fixed, and the product of the Weyl
translates of any denominator is an invariant denominator still nonzero
there.  This is also the central-parameter case of the prequotient
localization in \cite[Proposition~2.3]{DK25}.

Both the added loops and their cotangent partners have nonzero $H$-weights,
so $(D^+)^H$ sets all these coordinates to zero.  The enlarged moment map
then restricts to the original moment map.  Stability restricts to the
original stability as well: under a $G$-equivariant closed inclusion, the
Hilbert--Mumford test of a point and of every one-parameter subgroup is
unchanged.  Equivalently, semi-invariants lift by reductivity, and orbit
closures remain in the closed subvariety.  Hence
\[
 (D^+)^H=D,
\]
with $H$ acting trivially on the right.  Freeness of the stable actions now
gives a coefficient-compatible isomorphism
\begin{equation}\label{eq:loop-higgs-localization}
 H_H^*(X^+)\otimes_{R_0[z]}\Lambda
    \simeq (H^*(X)\otimes_\C\C[z])\otimes_{R_0[z]}\Lambda.
\end{equation}
The localization theorem used here does not require compactness of $D^+$
or $X^+$.

\emph{Descent.}
Let $I_B$ and $I_X$ be the kernels of the coefficient maps from $R_0$ to
$B_0$ and $H^*(X)$, respectively.  Equations
\eqref{eq:enlarged-flavor-isomorphism},
\eqref{eq:loop-coulomb-localization}, and
\eqref{eq:loop-higgs-localization} identify the localized quotients over the
same ring $\Lambda$.  Therefore
$I_B\Lambda=I_X\Lambda$.  Setting $z=1$ uses the identity
\[
 \Lambda/(z-1)\Lambda=(R_0)_{\mathfrak m_0}
\]
and yields
$(I_B)_{\mathfrak m_0}=(I_X)_{\mathfrak m_0}$.
Both ideals are homogeneous.  If $h\in I_B$ is homogeneous, there is
$s\in R_0$ with $s(0)\ne0$ and $sh\in I_X$.  The component of degree
$\deg h$ of this relation is $s(0)h\in I_X$, whence $h\in I_X$.
Interchanging the ideals proves $I_B=I_X$ globally.  This is equality of
ideals and gives the asserted graded isomorphism of quotient rings, with
all nilpotents retained.  The proof applies unchanged to empty stable loci; if there are
no nonzero gauge ranks, the assertion is immediate.
\end{proof}

\begin{theorem}\label{thm:full-torus}
For every datum in the statement of Theorem~\ref{thm:main}, the coefficient
map induces a canonical isomorphism of graded $R$-algebras
\[
 B_F\xrightarrow{\ \sim\ }H_F^*(X).
\]
\end{theorem}

\begin{proof}
At determinant stability, Proposition~\ref{prop:loop-removal} identifies
$B_0$ with ordinary cohomology, so $B_0$ is finite-dimensional.  For any
other regular stability, Lemma~\ref{lem:betti-independence} gives the same
total Betti number.  The graded coefficient-compatible surjection
$B_0\twoheadrightarrow H^*(X)$ of
Proposition~\ref{prop:kernel-inclusion} is therefore an isomorphism.
Lemma~\ref{lem:flavor-lift} proves the assertion for every commuting flavor
torus, with the generality supplied by Lemma~\ref{lem:trivial-summands}.
If a quotient is empty, the same argument gives $B_0=B_F=0$.
Uniqueness follows from surjectivity of the coefficient map $R\to B_F$.
\end{proof}
